\subsection{The polynomial sandwich}

The compact limit and the two exclusion estimates sandwich the full KKT event
between Poisson events whose score intervals differ by \(2\eta\).
\begin{theorem}
\label{thm:polynomial-stable-regime}
Fix $\alpha\in(0,1/4)$ and assume
\[
    \eps_n=n^{-\alpha+o(1)}.
\]
Let
\[
    \theta_\alpha=1-\sqrt\alpha.
\]
Let $S_+$ and $S_-$ be independent copies of the compensated Poisson process
\[
    S(\theta)
    =
    \int_\R e^{\theta x}\{N-\mu\}(dx),
    \qquad
    \mu(dx)=e^{-x}\,dx.
\]
Here \(S\) is interpreted using the half-line construction in
Proposition~\ref{prop:poisson-edge-field}.
For
\[
    0<\eta<\min\Bigl\{\theta_\alpha-\frac12,\sqrt\alpha\Bigr\},
\]
define
\[
    p_-(\alpha,\eta)
    =
    \mathbb P\Bigl\{
        \max_{\xi\in\{+,-\}}
        \sup_{\theta_\alpha-\eta\le\theta<1}
        S_\xi(\theta)<0
    \Bigr\},
\]
and
\[
    p_+(\alpha,\eta)
    =
    \mathbb P\Bigl\{
        \max_{\xi\in\{+,-\}}
        \sup_{\theta_\alpha+\eta\le\theta<1}
        S_\xi(\theta)\le0
    \Bigr\}.
\]
Then
\[
    p_-(\alpha,\eta)
    \le
    \liminf_{n\to\infty}
    \mathbb P\bigl\{|\supp(\wh G_n)|=1\bigr\}
    \le
    \limsup_{n\to\infty}
    \mathbb P\bigl\{|\supp(\wh G_n)|=1\bigr\}
    \le
    p_+(\alpha,\eta).
\]
Moreover
\[
    p_-(\alpha,\eta)>0,
    \qquad
    p_+(\alpha,\eta)<1.
\]
\end{theorem}

\begin{proof}
Put
\[
    \lambda_n
    =
    \frac{\eps_nu_n^2}{2(1-\eps_n)}.
\]
Then $\lambda_n=n^{-\alpha+o(1)}$, since $u_n^2/2=\log n+o(\log n)$ and $\eps_n\to0$.

The Poisson argument in Proposition~\ref{prop:poisson-edge-field} applies
jointly to both tails. After truncation to compact edge windows, the
multinomial rare-event formula gives independent Poisson limits; the
positive- and negative-edge truncation bounds in that proposition remove both
cutoffs. Thus the two compact fields converge
jointly to \((S_+,S_-)\).

Let
\[
    E_n=\{\sup_aD_n(a)>1\}.
\]
By Lemma~\ref{lem:one-atom-D}, the one-atom event is $E_n^c$.

For the upper bound on the one-atom probability, fix $\rho>0$ so small that
\[
    \theta_\alpha+\eta<1-\rho.
\]
On the compact interval $[\theta_\alpha+\eta,1-\rho]$ we have
\[
    (1-\theta)^2<\alpha
\]
uniformly. Hence
\[
    \lambda_n\theta^2/c_n(\theta)\to0
\]
uniformly there, and Proposition~\ref{prop:poisson-edge-field} plus Lemma~\ref{lem:uniform-centering} gives
\[
    \left\{
    \frac{D_{n,\xi}(\theta)-1}{c_n(\theta)}
    :
    \xi=\pm1,\ \theta\in[\theta_\alpha+\eta,1-\rho]
    \right\}
    \Rightarrow
    \{S_\xi(\theta):\xi=\pm1,\ \theta\in[\theta_\alpha+\eta,1-\rho]\}.
\]
Therefore
\[
    \liminf_n
    \mathbb P(E_n)
    \ge
    \mathbb P\Bigl\{
        \max_{\xi\in\{+,-\}}
        \sup_{\theta_\alpha+\eta\le\theta\le1-\rho}
        S_\xi(\theta)>0
    \Bigr\}.
\]
Letting $\rho\downarrow0$ and using $S_\xi(\theta)\to-\infty$ as $\theta\uparrow1$ gives
\[
    \limsup_n
    \mathbb P\bigl\{|\supp(\wh G_n)|=1\bigr\}
    \le
    p_+(\alpha,\eta).
\]

For the lower bound, Lemma~\ref{lem:subthreshold-envelope} excludes crossings
\(D_{n,\xi}(\theta)>1\) for
\[
    0\le\theta\le\theta_\alpha-\eta
\]
with high probability. Fix $R<\infty$. Lemma~\ref{lem:endpoint-exclusion} lets us discard
\[
    1-\rho\le\theta\le1+R/u_n^2
\]
at a cost that vanishes as $\rho\downarrow0$. For the part to the right of this window, use the event
\[
    \max_i Z_i\le u_n+\frac{R}{u_n},
    \qquad
    \max_i(-Z_i)\le u_n+\frac{R}{u_n}.
\]
Its complement has limiting probability $O(e^{-R})$, by Mills' ratio.
Intersect it with
\[
 |\bar Z_n|
 \le
 \frac{\eps_n}{2(1-\eps_n)}
 \left(u_n+\frac R{u_n}\right),
\]
an event whose probability tends to one because
\(\sqrt n\,\eps_nu_n\to\infty\). On this intersection, if
\(\theta>1+R/u_n^2\), then \(\xi a_n(\theta)\) lies strictly beyond the
corresponding centered sample extreme. Lemma~\ref{lem:one-atom-D} then makes
\(D_{n,\xi}(\theta)\) decreasing for larger \(\theta\), so any crossing there
would already force a crossing in the discarded endpoint window. We send
\(R\to\infty\) at the end. On the remaining compact interval
\[
    [\theta_\alpha-\eta,1-\rho],
\]
write
\[
    W_{n,\xi}(\theta)
    =
    \frac{D_{n,\xi}(\theta)-1}{c_n(\theta)}.
\]
The deterministic damping term in Proposition~\ref{prop:poisson-edge-field} is nonpositive, and the empirical-centering error is $o_{\mathbb P}(1)$ after division by $c_n(\theta)$ on this compact interval. Thus, for every fixed $\delta>0$, with probability tending to one, a crossing in the compact interval implies
\[
    \max_{\xi\in\{+,-\}}
    \sup_{\theta_\alpha-\eta\le\theta\le1-\rho}
    \frac{H_{n,\xi}(\theta)-1}{c_n(\theta)}
    \ge-\delta.
\]
The compact functional convergence and the portmanteau theorem for the closed
set $\{f:\sup f\ge-\delta\}$ bound the limsup of the compact crossing
probability by
\[
    \mathbb P\Bigl\{
        \max_{\xi\in\{+,-\}}
        \sup_{\theta_\alpha-\eta\le\theta\le1-\rho}
        S_\xi(\theta)\ge-\delta
    \Bigr\}.
\]
Letting \(\delta\downarrow0\) decreases these events to the event in which
both Poisson suprema are at most zero. Hence
\[
    \limsup_n
    \mathbb P(E_n)
    \le
    \mathbb P\Bigl\{
        \max_{\xi\in\{+,-\}}
        \sup_{\theta_\alpha-\eta\le\theta\le1-\rho}
        S_\xi(\theta)\ge0
    \Bigr\}
    +o_\rho(1)+O(e^{-R}).
\]
Letting first $\rho\downarrow0$ and then $R\to\infty$ gives
\[
    \liminf_n
    \mathbb P\bigl\{|\supp(\wh G_n)|=1\bigr\}
    \ge
    p_-(\alpha,\eta).
\]

We finish by proving that neither sandwich bound is trivial. For $p_->0$,
put $b=\theta_\alpha-\eta>1/2$. We show that one copy of $S$ is strictly
negative on $[b,1)$ with positive probability. Choose $M$ large. On the event
that $N$ has no points in $[-M,\infty)$, the contribution of this interval is
the deterministic negative function
\[
    -\int_{-M}^\infty e^{\theta x}e^{-x}\,dx
    =
    -\frac{e^{(1-\theta)M}}{1-\theta},
    \qquad b\le\theta<1.
\]
If $M>1/(1-b)$, the positive compensator $e^{(1-\theta)M}/(1-\theta)$ has minimum $eM$ over $[b,1)$. Let $T_M$ be the compensated contribution from $(-\infty,-M)$. By the $L^2$ isometry, applied also to the $\theta$-derivative, $\sup_{b\le\theta\le1}|T_M(\theta)|\to0$ in probability. Thus, for all large $M$, the event
\[
    \sup_{b\le\theta\le1}T_M(\theta)\le e\,M/2
\]
has positive probability. This event is independent of $\{N([-M,\infty))=0\}$, and their intersection has positive probability. On that intersection $S(\theta)<0$ for all $b\le\theta<1$. Since $S_+$ and $S_-$ are independent, both copies are negative on $[b,1)$ with positive probability. Hence $p_-(\alpha,\eta)>0$.

For $p_+<1$, choose any $\theta_0\in(\theta_\alpha+\eta,1)$.
The random variable $S_+(\theta_0)$ is centered, integrable, and nondegenerate.
Hence
$\mathbb P\{S_+(\theta_0)>0\}>0$, which gives $p_+(\alpha,\eta)<1$.
\end{proof}
\subsection{Continuity and strict monotonicity}

The sandwich identifies the interior limit, but it does not by itself show
that changing the exponent changes the limiting probability. We now prove the
interior regularity assertions in Proposition~\ref{prop:phase-curve-properties}.

\begin{proof}[Proof of the interior assertions in
Proposition~\ref{prop:phase-curve-properties}]
With
\[
    \mathcal M_b
    =
    \max_{\xi\in\{+,-\}}
    \sup_{b\le\theta<1}S_\xi(\theta),
\]
fix $0<r<1/4$ and write $\theta_r=1-\sqrt r$. The proof of $p_->0$ in Theorem~\ref{thm:polynomial-stable-regime}, applied with $b=\theta_r$, gives $\mathbb P\{\mathcal M_{\theta_r}<0\}>0$, and hence $p(r)>0$. The strict inequality $p(r)<1$ follows by evaluating one copy at any fixed $\theta_0\in(\theta_r,1)$: the centered variable $S_+(\theta_0)$ is nondegenerate, so it is positive with positive probability.

For monotonicity and continuity, put
\[
    F(b)=\mathbb P\{\mathcal M_b\le0\}.
\]
The random variables $\mathcal M_b$ decrease as $b$ increases, while $\theta_r=1-\sqrt r$ decreases as $r$ increases. Hence $p(r)=F(\theta_r)$ is nonincreasing in $r$.

To prove strictness, for one copy of the field put
\[
    \pi(b)=\mathbb P\Bigl\{\sup_{b\le\theta<1}S(\theta)\le0\Bigr\},
    \qquad b\in(1/2,1),
\]
so that \(F(b)=\pi(b)^2\). Fix \(1/2<b_0<b_1<1\). The event defining
\(\pi(b_0)\) is contained in the event defining \(\pi(b_1)\); to make the
containment strict, it suffices to find a positive-probability event on which
\(S\) is positive somewhere in \([b_0,b_1)\) but negative on \([b_1,1)\).

Under the change of variables \(y=e^{-x}\), the field has the representation
\[
    S(\theta)=\int_0^\infty y^{-\theta}\{\Pi(dy)-dy\},
\]
where \(\Pi\) is a unit-rate Poisson process. Choose a constant \(c\) strictly
between \(1/(1-b_0)\) and \(1/(1-b_1)\). For large \(Y\), set
\(k=\lfloor cY\rfloor\). Choose \(\delta>0\) small and put
\(L=(1+\delta)Y\). On the event \(E_Y\) that \(\Pi\) has no points in
\((0,Y)\) and exactly \(k\) points in \([Y,L]\), the contribution from
\((0,L]\), including its compensator, is bounded below at \(b_0\) by
\[
 kL^{-b_0}-\frac{L^{1-b_0}}{1-b_0}
 \ge mY^{1-b_0}
\]
and bounded above, uniformly for \(\theta\ge b_1\), by
\[
 kY^{-\theta}-\frac{L^{1-\theta}}{1-\theta}
 \le -mY^{1-\theta},
\]
for some \(m>0\) and all large \(Y\). These inequalities follow by first
setting \(\delta=0\), where they are precisely the two strict inequalities
defining \(c\), and then taking \(\delta\) small. Notice that \(E_Y\) has
positive probability for every fixed \(Y\).

It remains to control the independent compensated tail
\[
 R_L(\theta)=\int_L^\infty y^{-\theta}\{\Pi(dy)-dy\}.
\]
For \(Q_Y(\theta)=Y^{\theta-1}R_L(\theta)\), the \(L^2\) isometry gives,
uniformly on \([b_0,1]\),
\[
 \mathbb E|Q_Y(\theta)|^2+
 \mathbb E|Q_Y'(\theta)|^2\le \frac{C}{Y}.
\]
The derivative estimate follows by inserting
\(\log(Y/y)\) in the integrand. The Sobolev inequality therefore yields
\(\sup_{b_0\le\theta\le1}|Q_Y(\theta)|\to_{\mathbb P}0\). For all large
\(Y\), the tail preserves both strict signs above with positive probability.
By independence from \(E_Y\), the separating event has positive probability,
so \(\pi(b_0)<\pi(b_1)\), hence \(F(b_0)<F(b_1)\). Because
\(r\mapsto\theta_r\) is strictly decreasing, \(p(r)=F(\theta_r)\) is strictly
decreasing on \((0,1/4)\).

If $b_m\downarrow b$, then the domains $[b_m,1)$ increase and $\mathcal M_{b_m}\uparrow\mathcal M_b$ almost surely. Therefore $\mathbf 1_{\{\mathcal M_{b_m}\le0\}}\downarrow\mathbf 1_{\{\mathcal M_b\le0\}}$, and
\[
    \lim_mF(b_m)=F(b).
\]
If $b_m\uparrow b$, then the domains decrease and $\mathcal M_{b_m}\downarrow\mathcal M_b$ almost surely; the endpoint fact $S_\xi(\theta)\to-\infty$ as $\theta\uparrow1$ reduces the supremum to compact subintervals, where the paths are continuous. Hence $\mathbf 1_{\{\mathcal M_{b_m}\le0\}}\uparrow\mathbf 1_{\{\mathcal M_b<0\}}$. Lemma~\ref{lem:poisson-sup-no-atom} gives $\mathbb P\{\mathcal M_b=0\}=0$, so $F(b_m)\to F(b)$. This proves continuity of $F$, and hence of $p$, on the interior interval.
\end{proof}

Because the limiting supremum has no atom at zero, the strict and weak
no-crossing events have the same probability. The sandwich therefore yields
the interior case of Theorem~\ref{thm:exponent-phase-diagram}.
\begin{proposition}
\label{prop:interior-polynomial-limit}
Fix $\alpha\in(0,1/4)$ and assume
\[
    \eps_n=n^{-\alpha+o(1)}.
\]
Then
\[
    \mathbb P\bigl\{|\supp(\wh G_n)|=1\bigr\}
    \to
    p(\alpha).
\]
\end{proposition}

\begin{proof}
Write
\[
    \mathcal M_b
    =
    \max_{\xi\in\{+,-\}}
    \sup_{b\le\theta<1}S_\xi(\theta).
\]
Lemma~\ref{lem:poisson-sup-no-atom} gives $\mathbb P\{\mathcal M_{\theta_\alpha}=0\}=0$.
As $\eta\downarrow0$,
\[
    p_-(\alpha,\eta)
    =
    \mathbb P\{\mathcal M_{\theta_\alpha-\eta}<0\}
    \uparrow
    \mathbb P\{\mathcal M_{\theta_\alpha}<0\}.
\]
Here the increasing convergence follows from continuity on compact subintervals and the endpoint fact $S_\xi(\theta)\to-\infty$ as $\theta\uparrow1$.
Similarly,
\[
    p_+(\alpha,\eta)
    =
    \mathbb P\{\mathcal M_{\theta_\alpha+\eta}\le0\}
    \downarrow
    \mathbb P\{\mathcal M_{\theta_\alpha}\le0\}
    =
    \mathbb P\{\mathcal M_{\theta_\alpha}<0\}.
\]
The $\eta$-sandwich in Theorem~\ref{thm:polynomial-stable-regime} now forces the asserted limit.
\end{proof}

\begin{proof}[Proof of Theorem~\ref{thm:exponent-phase-diagram}]
For $0<r<1/4$, the assumption $r_n\to r$ is exactly $\eps_n=n^{-r+o(1)}$, so the result is Proposition~\ref{prop:interior-polynomial-limit}.

Suppose $r=0$. Fix $\alpha\in(0,1/4)$ and let $\eps_n^{(\alpha)}$ be chosen so that
$\eps_n^{(\alpha)}=n^{-\alpha}$.
Since $r_n\to0$, we have $\eps_n\ge\eps_n^{(\alpha)}$ for all large $n$. Coupling both samples through the same standardized variables $Z_i$, Lemma~\ref{lem:variance-monotonicity} gives
\[
    \bigl\{|\supp(\wh G_n^{(\eps_n^{(\alpha)})})|=1\bigr\}
    \subseteq
    \bigl\{|\supp(\wh G_n^{(\eps_n)})|=1\bigr\}
\]
for all large $n$, where $\wh G_n^{(\eps)}$ denotes the NPMLE computed from $\sqrt{1-\eps}\,Z_1,\dots,\sqrt{1-\eps}\,Z_n$.
Therefore
\[
    \liminf_{n\to\infty}
    \mathbb P\bigl\{|\supp(\wh G_n^{(\eps_n)})|=1\bigr\}
    \ge
    p(\alpha).
\]
Letting \(\alpha\downarrow0\) and using \(p(\alpha)\to1\) gives the limit one.

Finally suppose either $r_n\to r\ge1/4$, with $r=\infty$ allowed, or more generally $\liminf_n r_n\ge1/4$. Fix $\alpha\in(0,1/4)$. Then $\eps_n\le\eps_n^{(\alpha)}$ for all large $n$. Monotonicity gives
\[
    \bigl\{|\supp(\wh G_n^{(\eps_n)})|=1\bigr\}
    \subseteq
    \bigl\{|\supp(\wh G_n^{(\eps_n^{(\alpha)})})|=1\bigr\}
\]
for all large $n$. Hence
\[
    \limsup_{n\to\infty}
    \mathbb P\bigl\{|\supp(\wh G_n^{(\eps_n)})|=1\bigr\}
    \le
    p(\alpha).
\]
Letting \(\alpha\uparrow1/4\) and using \(p(\alpha)\to0\) gives the limit zero.
\end{proof}
